\section{2026-08-13 — Ratificación de XLM-T como modelo principal del clasificador}

Se ratificó \texttt{cardiffnlp/twitter-xlm-roberta-base} (XLM-T) como modelo principal del módulo de clasificación de lenguaje natural, en reemplazo de RoBERTuito, que pasa a ser la primera línea de comparación junto con BETO. El \textit{baseline} clásico de TF-IDF con clasificador lineal se mantiene sin cambios.

La decisión venía documentada como pendiente desde el diagnóstico del 8 de agosto y se resolvió al detectarse una contradicción concreta: el documento de arquitectura, redactado el 12 de agosto, ya nombraba a XLM-T, mientras que el resto del material seguía indicando RoBERTuito.

El argumento decisivo no fue el rendimiento sino la disponibilidad de datos. Las dos propiedades que el trabajo necesita estaban repartidas entre modelos distintos: RoBERTuito aporta la adaptación al registro de las redes sociales pero es monolingüe, y XLM-RoBERTa aporta la transferencia desde el inglés pero fue pre-entrenado sobre texto web de propósito general. Esa división importa porque los cerca de cuarenta mil ejemplos anotados de LIAR y FakeNewsNet están en inglés y solo resultan utilizables con un modelo multilingüe. Con un modelo monolingüe el conjunto de entrenamiento se habría reducido a FakeDeS —novecientos setenta y un ejemplos— más el corpus argentino todavía por construir, lo que trasladaba el riesgo completo a la etapa experimental. XLM-T reúne ambas propiedades porque parte de XLM-RoBERTa y continúa su pre-entrenamiento sobre alrededor de ciento noventa y ocho millones de publicaciones de Twitter en más de treinta idiomas.

La contrapartida queda declarada de manera explícita. Gouliev \textit{et al.} (2025) reportan una caída de entre ocho y doce puntos porcentuales de los modelos multilingües frente a los nativos de la lengua. Se asume esa contra sobre la base de que la medición se realizó con modelos multilingües sin adaptación de dominio, y de que el ajuste fino se ejecuta sobre datos en español; ninguno de los dos es un argumento concluyente, de modo que la comparación experimental contra RoBERTuito se conserva precisamente para verificar si la penalización aparece. Si el resultado contradice la elección, la elección se revierte.

Como consecuencia derivada se eliminó la estrategia de \textit{ensemble} de dos modelos que subsistía en el material de marco teórico. Esa recomendación proponía un modelo para publicaciones de redes y otro para artículos periodísticos, y había quedado sin sustento desde el acotamiento del alcance del 13 de junio: los medios de referencia son hoy fuente de evidencia y no objetos de clasificación, por lo que no hay artículos periodísticos que clasificar. Un único clasificador elimina además un servicio de inferencia del despliegue y su costo asociado.

Impacto en el documento: se incorporó la entrada bibliográfica \texttt{BarbieriEtAl2022}; se ampliaron la exposición de modelos y la tabla comparativa del Capítulo~2; se reformuló el pasaje que apoyaba la elección en Albtoush \textit{et al.} (2025), cuyo hallazgo se reinterpreta en términos de adaptación —de lengua o de dominio— frente a modelos de propósito general, para evitar una contradicción interna; se actualizó la tercera contribución declarada al cierre del Capítulo~2; y se sumó XLM-T a la enumeración de arquitecturas candidatas de los objetivos específicos del Capítulo~1, que conserva su redacción agnóstica.
